%============================================================
% HALAMAN KATA PENGANTAR
%============================================================
\cleardoublepage
\chapter*{KATA PENGANTAR}
\addcontentsline{toc}{chapter}{KATA PENGANTAR}
\thispagestyle{fancy}

Puji syukur penulis panjatkan ke hadirat Allah SWT atas segala rahmat dan karunia-Nya sehingga proposal Proyek Akhir ini dapat diselesaikan dengan baik. Proposal dengan judul \textit{``Implementasi Agentic DevOps Monitoring Agents Berbasis LLM untuk Meningkatkan Observability dan Self-Healing Backend FoodLAB''} disusun sebagai syarat mengikuti seminar proposal Proyek Akhir di Politeknik Elektronika Negeri Surabaya (PENS).

Penulis menyadari bahwa penyelesaian proposal ini tidak terlepas dari dukungan, bimbingan, dan kontribusi berbagai pihak. Oleh karena itu, penulis menyampaikan terima kasih yang sebesar-besarnya kepada:

\begin{enumerate}[label=\arabic*., leftmargin=*]
  \item Bapak Adam Shidqul Aziz, S.Kom., M.Kom., selaku dosen pembimbing I yang telah memberikan arahan, masukan teknis, dan motivasi selama proses penyusunan proposal ini.
  \item Ibu Umi Sa'adah, S.Kom., M.Kom., selaku dosen pembimbing II yang telah memberikan bimbingan metodologis dan dukungan penuh dalam penyelesaian proposal ini.
  \item Bapak Willy Achmat Fauzi beserta seluruh pihak PT Sinergi Dimensi Informatika yang telah memberikan kesempatan dan dukungan dalam konteks penerapan sistem \textit{monitoring} pada lingkungan nyata.
  \item Seluruh dosen dan staf Program Studi D4 Teknik Informatika PENS atas fasilitas dan lingkungan akademik yang kondusif.
  \item Rekan-rekan mahasiswa atas diskusi, masukan, dan semangat kebersamaan selama pengerjaan proposal ini.
\end{enumerate}

Penulis menyadari bahwa proposal ini masih memiliki kekurangan dan keterbatasan. Masukan serta saran yang membangun dari berbagai pihak sangat diharapkan untuk penyempurnaan pada tahap pelaksanaan dan pelaporan Proyek Akhir selanjutnya.

\vspace{1.5cm}
\begin{flushright}
  Surabaya, Juli 2026\\[0.5cm]
  Penulis
\end{flushright}